\chapter{Troubleshooting}
\label{chap:diagnostico}

When a fault occurs, check the system in the same order as bring-up: PCIe connection,
driver, Python environment, model and data. \Cref{tab:troubleshooting} summarizes
common faults and where to start looking.

\begin{longtable}{@{}L{0.25\textwidth}L{0.34\textwidth}L{0.33\textwidth}@{}}
\caption{Quick troubleshooting guide.}
\label{tab:troubleshooting}\\
\toprule
\textbf{Symptom} & \textbf{Check} & \textbf{Action} \\
\midrule
\endfirsthead

\toprule
\textbf{Symptom} & \textbf{Check} & \textbf{Action} \\
\midrule
\endhead

The board does not appear in \code{lspci}. &
Power, FFC cable, assembly and Raspberry Pi 5 PCIe configuration. &
Check the physical connection and PCIe configuration first. \\

The board appears on PCIe, but \code{/dev/akida*} does not exist. &
Kernel version, headers, \code{dkms status}, \code{lsmod} and \code{dmesg} messages. &
Check driver installation and compatibility. \\

\code{/dev/akida*} exists, but \code{akida.devices()} is empty. &
Active virtual environment, Akida version and device permissions. &
Check that you are using the right environment and that the user can access the device. \\

\code{import akida} fails. &
Python version and packages installed in the active environment. &
Activate the right virtual environment or reinstall Akida in it. \\

The FBZ file does not load. &
File path, Akida version and model compatibility. &
First check the test model included with the repository. \\

The model cannot map with \code{hw\_only=True}. &
\code{model.summary()} output, layers used and model version. &
Check that the model supports execution entirely on the AKD1000. \\

The input is rejected. &
Data shape, type, range and order. &
Compare the input with \code{model.input\_shape} and the preprocessing used during conversion. \\

Execution finishes, but the output is not as expected. &
Preprocessing, sample order and use of \code{forward()} or \code{predict()}. &
First test an input with a known expected result. \\

Latency varies widely between runs. &
Clock mode, mapping mode, warm-up, host load and number of repetitions. &
Repeat the measurement with a fixed configuration. \\

No power measurement appears. &
Power measurement enabled and statistics returned by Akida. &
Check the configuration and increase the test duration if needed. \\

\bottomrule
\end{longtable}

\section{Information for diagnosing a fault}

The repository's tool can collect much of the information needed to locate a fault:

\begin{lstlisting}[language=bash]
mkdir -p evidence
akd1000-bringup doctor --json > evidence/doctor.json
\end{lstlisting}

For more detail, also save the system and driver status:

\begin{lstlisting}[language=bash]
uname -a > evidence/uname.txt
lspci -nnk > evidence/lspci-nnk.txt
dkms status > evidence/dkms.txt
ls -l /dev/akida* > evidence/device-nodes.txt 2>&1
python -m pip freeze > evidence/pip-freeze.txt
sudo dmesg | grep -iE 'akida|edma|pcie' > evidence/driver-log.txt
\end{lstlisting}

When reporting a fault, also retain the command that caused it and the full error
message. This usually helps distinguish a PCIe, driver, Python environment, model
or data problem.
